\section{Decision-scale precision and the cost of a qualified NBO}\label{sec:precision96}
An approximation can save continuation queries and still leave an economic comparison unresolved. The investment experiment in Section~\ref{sec:study95} illustrates this distinction: the original policy intervals establish equivalence only at a margin equal to four fifths of the maximal first-date investment charge. We now connect numerical integration, policy admission, statistical resolution and the cost of deployment. The original economic law and the NBO construction remain the benchmark. Neither a smaller numerical interval nor a new admission rule is silently attributed to the unchanged historical actor.

\subsection{Risk quadrature from the convexity of the original value}
Fix a current state and feasible action. Let $v(z)=V_1(F(x,a,z))$ on $[-R,R]$, where $R=1/32$. Proposition~\ref{prop:curvature95} implies that $v$ is convex and $15$-Lipschitz. This conclusion concerns the true one-date value. The complete terminal controller need not be convex: its original bound $V_1\leq V_1^{\widehat\pi}\leq V_1+\varepsilon_1$, $\varepsilon_1=2^{-20}$, will be added separately.

Let $z_0,\ldots,z_n$ be the endpoints of equal innovation bins, and let $m_i=(z_{i-1}+z_i)/2$. Define $M_n$ to put probability $1/n$ at each midpoint. Define $E_n$ to put probability $1/(2n)$ at each of the two endpoints of every bin, counting an interior endpoint twice. These are numerical measures used for enclosures, not alternative economic transition laws.

\begin{proposition}[Convex-order risk enclosure]\label{prop:quadrature96}
For a convex $K$-Lipschitz function $v$ and the uniform law $P$ on $[-R,R]$,
\begin{equation}\label{eq:quad96}
 \rho_\alpha^{M_n}(v)\leq\rho_\alpha^P(v)
                 \leq\rho_\alpha^{E_n}(v),\qquad
 \rho_\alpha^{E_n}(v)-\rho_\alpha^{M_n}(v)
          \leq\frac{KR}{(1-\alpha)n^2}.
\end{equation}
If midpoint lower values and endpoint upper values have errors at most $e_-$ and $e_+$ respectively, the computed width is at most the last term plus $e_-+e_+$. A terminal controller uniformly within $\varepsilon_1$ of $V_1$ adds only $\beta\varepsilon_1$ to its first-date upper cost. All numerical endpoint arithmetic must be enclosed separately.
\end{proposition}

The proof combines the established convex midpoint--trapezoid order with the AVaR threshold representation; neither underlying principle is claimed as new \citep{alomari2016hermite,RockafellarUryasev2000}. Their use here changes the implementable original-law precision account. The previous full-bin Lipschitz allowance decays as $n^{-1}$. The new width decays as $n^{-2}$ without assuming a differentiable risk-averse value or a fixed maximizing risk density. Both numerical risks are finite weighted tail calculations. Ignoring the half weights at the two outer endpoints would compute a different bound.

\begin{corollary}[Sharpness on an explicit neural continuation]\label{cor:sharpness96}
Let $\alpha\in\{1/2,3/4\}$ and let $n$ be a positive multiple of eight. For $K,R>0$, put $c_n=R/n$ and
\[
 v_n(z)=K\{\sigma(z-c_n)+\sigma(c_n-z)\}=K|z-c_n|.
\]
This convex two-unit ReLU continuation satisfies
\[
 \rho_\alpha^{E_n}(v_n)-\rho_\alpha^{M_n}(v_n)
       =\frac{KR}{(1-\alpha)n^2}.
\]
Consequently neither the exponent nor the leading constant in the noiseless width bound of Proposition~\ref{prop:quadrature96} can be reduced uniformly over convex $K$-Lipschitz continuations, even on this explicit neural subclass. The conclusion concerns the stated midpoint and endpoint rules, not a lower bound for all integration algorithms or a superiority claim for neural representations.
\end{corollary}


The same argument applies to a rectangular product innovation law when $v$ is convex in each coordinate. If coordinate $j$ has radius $R_j$, Lipschitz constant $K_j$ and $n_j$ bins, the product midpoint and product endpoint rules have width at most
\[
 \frac{1}{1-\alpha}\sum_j\frac{K_jR_j}{n_j^2}.
\]
The number of evaluation points still grows as a product of the resolutions. This extension therefore retains the two-innovation cost rather than claiming that the curse of dimension has disappeared. It does not automatically apply to the possibly nonconvex minimum-of-cones representation of the general constructive backend. Convexity of a supplied continuation, or of the original value being enclosed, is a substantive hypothesis.

For the investment model, a midpoint affine support also supplies a feasible density witness for the all-original-action lower quadratic in Theorem~\ref{thm:continuous95}. The new endpoint upper value and that lower quadratic therefore give
\[
 \underline V_0(x)\leq V_0(x)\leq J_0^{\widehat\pi}(x)
                   \leq\overline J_0(x).
\]
The policy upper and lower costs use the same continuous innovation law. The state-cover residual remains a whole-domain statement, whereas an admission inequality is checked at the queried state. These two quantifiers are not interchangeable.

\subsection{Admission changes the statistical support}
A stored continuation first proposes an action. A qualified service accepts that action only when its complete-policy bound is at most $\tau$ against the original optimum. Otherwise it pays for the declared conventional action and verifies that action. An exact finite-cover fallback makes this a complete measurable service on the declared dyadic initial domain, rather than a controller defined only on successful sample rows. The fallback is part of the policy, and a rejected proposed action remains in the record.

\begin{theorem}[Qualification and simultaneous resolving power]\label{thm:power96}
Consider $A$ complete services, fixed independently of an evaluation sample, each satisfying
$0\leq J^a(x)-V^*(x)\leq\tau$ for every state in the declared initial domain. Suppose each reported policy-cost enclosure has width at most $w$. For $N$ independent initial states, let $[d_i^-,d_i^+]$ be the difference of two cost enclosures, intersected with $[-\tau,\tau]$. For any $r>0$, all $H$ declared pairwise intervals
\begin{equation}\label{eq:paired96}
 \left[\frac1N\sum_i d_i^- -r,
       \frac1N\sum_i d_i^+ +r\right]
\end{equation}
cover their original-law mean contrasts with probability at least
$1-2H\exp\{-Nr^2/(2\tau^2)\}$. Their half-widths are at most $r+w$, apart from the separately enclosed arithmetic used to report endpoints.
\end{theorem}

The global service definition is essential. Success on the observed states alone does not establish the support used in this theorem. Nor may a failed fitted proposal be deleted from a nominally unchanged-policy sample. The present service records its proposal, acceptance decision, final action, numerical refinement and exact-fallback entry. The original R95 policy experiment remains a different, unchanged experiment.

We use $\tau=1/256$, $w=1/8192$, $N=4096$, $H=12$ across two risk levels, and $r=1/4096$. Then $2H\exp(-8)<0.05$, while
\begin{equation}\label{eq:power96}
 r+w=\frac3{8192}<\frac1{2560}=\frac{\epsilon}{2},
 \qquad \epsilon=\frac1{1280}=\frac{h(1/4)}{100}.
\end{equation}
Thus the sample size and numerical precision are fixed together to resolve an average policy-cost difference at one percent of the largest first-date investment charge. The margin is eighty times smaller than the earlier $1/16$ margin. It is an explicit model-relative purchaser tolerance, not an estimated welfare weight or external monetary calibration. We report which comparisons actually meet it; the design alone implies no favorable sign.

\subsection{Acquisition, reuse and an economic purchase}
Let $K_a$ be a one-time acquisition charge for service $a$, and $c_a$ its cost per future workload. The acquisition account includes labels, feature search, fitting, validation, serialization and the chosen audit campaign. A fresh-process workload includes loading, identity checks, action proposal, admission, fallback, numerical certification, output and shutdown. Let a workload contain $B$ initial-state decisions and let $d=\mathbb E[J^a-J^b]$ be the per-decision original policy-cost contrast, $\lambda\geq0$ the price of computational work in the same normalized cost units, and $F$ an additional fixed fee difference. For $m$ future workloads, the total normalized difference is
\begin{equation}\label{eq:purchase96}
 m\{B d+\lambda(c_a-c_b)\}+\lambda(K_a-K_b)+F.
\end{equation}
For the recorded timing workloads $B=32$; for a single decision $B=1$. This expression distinguishes a policy-quality loss from the value of saved computation. A price is not inferred from seconds, a coefficient count, or an integration bin.

\begin{proposition}[Robust acquisition-and-reuse region]\label{prop:purchase96}
Suppose the policy contrast has simultaneous upper bound $d^+$, the workload saving $c_b-c_a$ has a valid lower bound $s^-$, and the incremental acquisition and fee have upper bounds $k^+$ and $F^+$. Then service $a$ weakly improves total normalized cost whenever
\begin{equation}\label{eq:adopt96}
 m(\lambda s^- -B d^+)\geq\lambda k^++F^+.
\end{equation}
If $\lambda s^- -B d^+>0$ and the right side is nonnegative, the smallest sufficient integer workload count is its ratio rounded upward. If the denominator is nonpositive and the right side is positive, no positive workload count is certified by these bounds. Separately, a purchaser's quality condition is $d^+\leq\epsilon$; it must not be inferred from a positive computational saving.
\end{proposition}

The timing repetitions below produce descriptive machine-specific ranges, not probabilistic guarantees for arbitrary future machines. Inserting their extrema in \eqref{eq:adopt96} gives a sensitivity calculation conditional on that timing envelope. It does not turn five repeated processes into a confidence bound on long-run prices or runtime. The theorem states the economic decision that adequately measured cost bounds would license, while the experimental record states precisely which costs were observed.
